% page 164 of the facsimile (image p-163 of the scan)
% block 188.0 x 259.3 mm ; left margin 8.0 mm
\pgc{-12.700mm}{0.000mm}{
\l{}
\l{}
\l{}
\l{\cel{15}154.}
\l{}
\l{}
\l{\cel{12}\sou{extradoso}. (\sou{arkitekt.)}\cel{2}Surfaco konvexa quan vulto prizentas}
\l{\cel{15}ad-extere (opoze a \textquotedbl{}intradoso\textquotedbl{}).}
\l{}
\l{\cel{12}\sou{extraktar}. (\sou{trans.}) I. Retrotirar (kozo) de loko ube lu esas}
\l{\cel{15}sinkita, stekita, nature od acidente. - II. Tirar de libro,}
\l{\cel{15}de manuskripto (un o plura texto-peci quin onu kopias).-}
\l{\cel{15}III. Deprenar de substanco, per deskompozar lu (elemento}
\l{\cel{15}quan onu separas de lu).- IV. (\sou{aritm.}) Extraktar radiko qua}
\l{\cel{15}dratala, kubala : ektirar ta radiko de ta nombro per des-}
\l{\cel{15}kompozar ica. - V. (\sou{kemio}). Extraktar salo : per analizo}
\l{\cel{15}kemiala. - DEFIRS.}
\l{}
\l{\cel{12}\sou{extraordinara}. I. Qua ne resortisas a la regulo, a la normo.}
\l{\cel{15}- II. Qua diferas de to quon onu vidas ordinare. - III. Qua}
\l{\cel{15}esas super la nivelo ordinara di la kozi, di la personi.-}
\l{\cel{15}DEF.}
\l{}
\l{\cel{12}\sou{extravagar}. (\sou{netran}s)Agar, dicar kozi qui ne korespondas a la}
\l{\cel{15}raciono komuna. - DEFIS.}
\l{}
\l{\cel{12}\sou{extrema}.- I. Qua esas tote ye la fino. - II. Qua esas tote ye}
\l{\cel{15}la lasta grado. - III.Qua esas ye grado tre alta. - IV.}
\l{\cel{15}Qua agas til la lasta punto. - DEFIS.}
\l{}
\l{\cel{12}\sou{extrinseka(filoz.}) Qua tiresas, cherpesas, ne de la kozo ipsa,}
\l{\cel{15}ma de cirkonstanci acesora, extera. - EFIS.}
\l{}
\l{\cel{12}\sou{exudar}. (\sou{netrans.}) Ekfluar preske nesenteble, quale sudoro.}
\l{\cel{15}- DEFIS.}
\l{}
\l{\cel{12}\sou{exulcero.}\cel{1}(\sou{medic.)}\cel{2}Ulcero neprofunda.}
\l{}
\l{\cel{12}\sou{exultar}. (\sou{netrans.}) Kaptesar da joyo profunda o da sentimento}
\l{\cel{15}(amo, amoro) qua su expresas per movi, kanti, e c. - EFI.}
\l{}
\l{\cel{12}\sou{exutorio}.\cel{2}I. (\sou{medic.})\cel{2}Ulcero artificala, destinita a plear}
\l{\cel{15}la rolo di deturnivo, por produktar ed entratenar pusifo}
\l{\cel{15}lokala. - II. Moyeno uzata por eskartar ulo mala, por des-}
\l{\cel{15}embarasar su ye lu. - FIS.}
\l{}
\l{\cel{12}\sou{-eyo}-. Sufixo qua signifikas \textquotedbl{}loko destinita por\textquotedbl{}.}
\l{}
\l{\cel{12}\sou{-ez-}. (\sou{gram.})\cel{3}Dezinenco di la imperativo.}
\l{}
\l{\cel{12}\sou{ezofago}. (\sou{anat.}) Tubo membrana qua iras de la faringo an la}
\l{\cel{15}stomako,e duktas ad ica la nutrivi. - EFIS.}
\l{}
\l{}
\l{}
\l{}
\l{}
\l{}
\l{}
\l{\cel{43}---}
}
